\documentclass[12pt,a4paper]{article}
% Сборка из корня репозитория:
% pdflatex -interaction=nonstopmode -halt-on-error -output-directory=out Tex/lecture-06.tex
\usepackage[utf8]{inputenc}
\usepackage[T2A]{fontenc}
\usepackage[russian]{babel}
\usepackage{amsmath,amssymb,amsthm,mathtools}
\usepackage{geometry}\geometry{margin=2.2cm}
\usepackage{microtype,parskip,enumitem,tikz}
\usetikzlibrary{arrows.meta,positioning,calc}
\newcommand{\Z}{\mathbb Z}
\newcommand{\Q}{\mathbb Q}
\newcommand{\R}{\mathbb R}
\newcommand{\CC}{\mathbb C}
\newcommand{\F}{\mathbb F}
\newcommand{\PP}{\mathbb P}
\newcommand{\tors}{\mathrm{tors}}
\DeclareMathOperator{\ord}{ord}
\DeclareMathOperator{\charop}{char}
\newcommand{\leg}[2]{\left(\frac{#1}{#2}\right)}
\setlength{\emergencystretch}{1.5em}
\allowdisplaybreaks[1]
\raggedbottom

% Все рисунки построены функциями; растровые фрагменты не используются.
% Односвязная модель: x=t^2, y=t*sqrt(t^4+1), то есть y^2=x^3+x.
% Двухкомпонентная модель: y^2=0.3(x+3)(x+1)(x-1).
% Овал: x=-2-cos(t), y=sqrt(.3)*sin(t)*sqrt(3+cos(t)).
% Неограниченная компонента: x=1+t^2,
% y=sqrt(.3)*t*sqrt((t^2+4)(t^2+2)).
\pgfmathdeclarefunction{onex}{1}{\pgfmathparse{#1*#1}}
\pgfmathdeclarefunction{oney}{1}{\pgfmathparse{#1*sqrt(#1*#1*#1*#1+1)}}
\pgfmathdeclarefunction{ovalx}{1}{\pgfmathparse{-2-cos(#1)}}
\pgfmathdeclarefunction{ovaly}{1}{\pgfmathparse{sqrt(.3)*sin(#1)*sqrt(3+cos(#1))}}
\pgfmathdeclarefunction{branchx}{1}{\pgfmathparse{1+#1*#1}}
\pgfmathdeclarefunction{branchy}{1}{\pgfmathparse{sqrt(.3)*#1*sqrt((#1*#1+4)*(#1*#1+2))}}
\newcommand{\ECoval}[1]{%
  \draw[#1,domain=0:360,samples=145,smooth,variable=\t]
    plot ({ovalx(\t)},{ovaly(\t)});}
\newcommand{\ECbranch}[2]{%
  \draw[#1,domain=-#2:#2,samples=145,smooth,variable=\t]
    plot ({branchx(\t)},{branchy(\t)});}
\tikzset{curve/.style={line width=.85pt},auxline/.style={line width=.65pt},
  axes/.style={->,line width=.5pt},every picture/.style={>=Latex}}

% Точки на двухкомпонентной модели при y=+-0.5.
\def\xa{-2.886326}\def\xb{-1.210671}\def\xc{1.096996}
\def\xr{2.284158}\def\yr{2.585650}\def\slope{.596782}

\newcommand{\ProjectionSketch}{%
\begin{tikzpicture}[x=1.4cm,y=.8cm,font=\small]
  \draw[axes] (-1.15,0)--(2.2,0);
  \draw[curve,domain=-1.38:1.38,samples=145,smooth,variable=\t]
    plot ({onex(\t)},{oney(\t)});
  \draw[auxline] (-.55,-1.39)--(1.94,2.98);
  \draw[auxline] (-.44,.45)--(.99,-1.57);
  \fill (1.5,2.20794) circle (1.6pt) node[above left] {$P$};
  \fill (.080128,-.283978) circle (1.6pt) node[right=5mm] {$Q$};
\end{tikzpicture}}

\newcommand{\RootPictures}{%
\begin{center}
\begin{minipage}{.44\textwidth}\centering
\begin{tikzpicture}[x=1.15cm,y=.72cm,font=\small]
  \draw[axes] (-1.15,0)--(2.65,0);
  \draw[axes] (0,-2.7)--(0,2.7);
  \draw[curve,domain=-1.3:1.3,samples=145,smooth,variable=\t]
    plot ({.55+onex(\t)},{oney(\t)});
  \node[anchor=west] at (1.2,-.65) {1 корень};
\end{tikzpicture}
\end{minipage}\hfill
\begin{minipage}{.54\textwidth}\centering
\begin{tikzpicture}[x=.94cm,y=.72cm,font=\small]
  \draw[axes] (-3.5,0)--(2.9,0);
  \draw[axes] (0,-2.7)--(0,2.7);
  \ECoval{curve}\ECbranch{curve}{1.15}
  \node[anchor=west] at (1.55,.4) {3 корня};
  \node[anchor=west] at (2.1,-1.6) {$E(\R)$};
\end{tikzpicture}
\end{minipage}
\end{center}}

\newcommand{\ChordPicture}{%
\begin{tikzpicture}[x=1.35cm,y=.9cm,font=\small]
  \draw[axes] (-3.5,0)--(3.55,0);
  \draw[axes] (0,-2.2)--(0,2.3);
  \ECoval{curve}\ECbranch{curve}{1.05}
  % Горизонтальные масштабы и положение точек повторяют рукопись.
  % Секущая проходит через две точки овала и правую компоненту.
  \draw[auxline] (-3.35,{.28*(-3.35)+.54})--(2.8,{.28*2.8+.54});
  \fill (-2.964011,-.289923) circle (1.7pt) node[below left] {$P$};
  \fill (-1.050414,.245884) circle (1.7pt) node[above left=7pt] {$Q$};
  \fill (1.275759,.897213) circle (1.7pt);
  \node[anchor=west] at (1.8,.55) {$PQ\ (x,y)$};
  \draw[dashed] (1.275759,.897213)--(1.275759,-.897213);
  \fill (1.275759,-.897213) circle (1.7pt);
  \node[anchor=west,align=left] at (1.75,-.8) {$(x,-y)$\\$P+Q$};
\end{tikzpicture}}

\newcommand{\AssociativityStart}{%
\begin{tikzpicture}[x=1.35cm,y=.85cm,font=\small]
  \draw[axes] (-3.45,0)--(3.55,0);
  \draw[axes] (0,-1.65)--(0,1.65);
  \ECoval{curve,red}
  \fill (\xa,-.5) circle (1.7pt) node[below left] {$P_1$};
  \fill (\xb,.5) circle (1.7pt) node[above right] {$P_2$};
  \fill (\xa,.5) circle (1.7pt) node[above left] {$P_3$};
\end{tikzpicture}}

% #1: цвет овала; #2: подпись нижней правой точки (как в рукописи).
\newcommand{\AssociativityConstruction}[2]{%
  \draw[axes] (-3.55,0)--(3.65,0);
  \draw[axes] (0,-3.25)--(0,3.25);
  \ECoval{curve,#1}\ECbranch{curve}{1.25}
  \draw[auxline] (-3.45,{-.5+\slope*(-3.45-\xa)})--(2.58,{-.5+\slope*(2.58-\xa)})
    node[pos=.76,above,sloped] {$L_1$};
  \draw[auxline] (\xa,.5)--(\xr,-\yr) node[pos=.78,below,sloped] {$L_2$};
  \draw[auxline] (\xa,.5)--(\xr,.5) node[pos=.5,above] {$\ell_1$};
  \draw[auxline] (\xa,-.5)--(\xc,-.5) node[pos=.78,above] {$\ell_2$};
  \draw[dashed] (\xc,-1.25)--(\xc,1.95) node[pos=.94,right] {$L_3$};
  \draw[dashed] (\xr,-2.92)--(\xr,2.95) node[pos=.56,right] {$\ell_3$};
  \fill (\xa,-.5) circle (1.7pt) node[below left] {$P_1$};
  \fill (\xa,.5) circle (1.7pt) node[above left] {$P_3$};
  \fill (\xb,.5) circle (1.7pt) node[above right=6pt] {$P_2$};
  \fill (\xb,-.5) circle (1.7pt);
  \node[anchor=north] at (\xb,-1.15) {$P_1(P_2+P_3)$};
  \fill (\xc,.5) circle (1.7pt);
  \node[anchor=west] at (1.35,.75) {$P_2P_3$};
  \fill (\xc,-.5) circle (1.7pt);
  \fill (\xr,\yr) circle (1.7pt) node[right] {$P_1P_2$};
  \fill (\xr,-\yr) circle (1.7pt) node[right] {$#2$};
}

\begin{document}
\section*{\centering Лекция 6}

% Рукопись, с. 1. Все утверждения и исторические пометы сохранены.
\section{Род кривой и рациональные точки}

\textbf{Memento.} Пусть $C$ --- неприводимая кривая, $\deg C=n$. Тогда
\[
g(C)=\frac{(n-1)(n-2)}2-m,
\qquad m\text{ --- число двойных точек}.
\]

\textbf{Теорема.} Пусть $E=\{y^2=x^3+Ax+B\}$ над $k\subset\CC$. Тогда
\begin{align*}
E\text{ гладкая}
&\Longleftrightarrow g(E)=1
\Longleftrightarrow \Delta=-4A^3-27B^2\ne0\\
&\Longleftrightarrow \text{все корни }x^3+Ax+B\text{ простые},
\end{align*}
где $\Delta$ --- дискриминант $f(x)=x^3+Ax+B$.

\textbf{Факт (Гильберт, Гурвиц).} Если $g(C)=0$, $\deg C=n>3$, то
$C$ бирациональна кривой степени $n-2$.

\textbf{Мораль.} Такие кривые бирациональны прямой или конике.

\textbf{Факт (Пуанкаре).} Любая кривая $C$ с рациональной точкой и
$g(C)=1$ бирациональна кубике.

\textbf{Теорема (гипотеза Морделла, теорема Фалтингса).}
Если $g(C)>1$, то
\[
|C(\Q)|<\infty,
\qquad C(\Q)\text{ --- рациональные точки}.
\]

\subsection*{Вопрос с кривыми рода 1}
Пусть $E$ --- кривая рода $1$ (кубика). По теореме Морделла
\[
E(\Q)\simeq\Z^r\oplus
\underbrace{\Z/p_1^{s_1}\Z\oplus\dots\oplus\Z/p_m^{s_m}\Z}
_{E(\Q)_{\tors}\text{ --- кручение}}.
\]

\textbf{Гипотеза.} Ограничено ли $r$? (Noam Elkies, 2024, $r=29$.)

\textbf{Теорема (Mazur, 1976).}
\[
|E(\Q)_{\tors}|\in\{1,\dots,10,12,16\}.
\]
Все значения достигаются.

$\Q$ можно заменить на любое числовое поле $k$.

\textbf{Теорема.} Любая гладкая кубика с рациональной точкой бирациональна
кубике
\[
y^2=f(x),\qquad \deg f=3,\qquad\Delta(f)\ne0.
\]

\newpage
% Рукопись, с. 2.
\section{Приведение кубики к виду $y^2=f(x)$}

\begin{minipage}[t]{.58\textwidth}
\vspace{0pt}
Пусть $g(x,y)$ --- уравнение исходной кривой,
\[
P\in C=\{g(x,y)=0\}
\]
--- рациональная точка. Можем сделать $Q$ началом координат.
\end{minipage}\hfill
\begin{minipage}[t]{.4\textwidth}\vspace{-4mm}\centering
\ProjectionSketch
\end{minipage}

\[
g(x,y)=\varphi_3(x,y)+2\varphi_2(x,y)+\varphi_1(x,y),
\]
где $\varphi_i(x,y)$ --- многочлен однородной степени $i$.

Подставим $L=\{y=tx\}$ в уравнение:
\[
\psi_3x^2+2\psi_2x+\psi_1=0,
\qquad \psi_i(t)=\varphi_i(1,t),
\qquad \varphi_i(x,tx)=x^i\varphi_i(1,t).
\]
\[
x=\frac{-\psi_2\pm\sqrt{\psi_2^2-\psi_3\psi_1}}{\psi_3}.
\]
Следовательно,
\[
(x\psi_3+\psi_2)^2=\psi_2^2-\psi_3\psi_1,
\qquad \deg(\psi_2^2-\psi_3\psi_1)\le4.
\]

\textbf{1) Если $\deg(\psi_2^2-\psi_3\psi_1)<4$.} Обозначим
\[
f=\psi_2^2-\psi_3\psi_1,\qquad
X=t=\frac yx,\qquad Y=\psi_3x+\psi_2.
\]
Тогда $Y^2=f(X)$. При этом $\deg f=3$, так как $g=1$, а род ---
бирациональный инвариант.

\textbf{2) Если степень равна $4$.} Существует $y=t_0x$, для которой
\[
\Delta(t)=(t-t_0)h(t),\qquad \deg h=3,
\qquad \Delta(t)=\psi_2^2(t)-\psi_3(t)\psi_1(t).
\]
Положим $t=t_0+\dfrac1X$. Тогда
\[
\Delta\!\left(t_0+\frac1X\right)
=\frac1Xh\!\left(t_0+\frac1X\right)
=\frac{X^3h\!\left(t_0+\frac1X\right)}{X^4}.
\]
\[
Y=X^2(\psi_3x+\psi_2),\qquad \Delta(f)\ne0,
\quad\text{так как }g=1.
\]

\textbf{Следствие.} Любую гладкую кубику с рациональной точкой можно
привести к виду
\[
y^2=x^3+Ax+B,\qquad 4A^3-27B^2\ne0.
\]

\textbf{Доказательство.} По замене получаем $y^2=ax^3+bx^2+cx+d$.
В рукописи записаны замены
\[
x\longmapsto a^2x\ \Longrightarrow\ a=1;\qquad
x\longmapsto x-\frac b3\quad\text{--- убрать квадратный член}.
\]

\newpage
% Рукопись, с. 3.
\section{Эллиптические кривые}

\textbf{Определение.} Эллиптическая кривая над $k$ --- проективная кривая
\[
Y^2Z=X^3+AXZ^2+BZ^3,\qquad
\Delta=-4A^3-27B^2\ne0,
\]
где $\charop(k)\ne2,3$.

\textbf{Замечание.} $g(E)=1$ и $E$ гладкая.
\[
E=\{y^2=x^3+Ax+B\}\cup\{O\},\qquad O=(0:1:0).
\]

\RootPictures

\section{Групповой закон}
Пусть $P,Q\in E(k)$. Обозначим через $PQ$ третью точку пересечения
прямой $L_{P,Q}$ с $E$:
\[
PQ=L_{P,Q}\cap E(k)\quad\text{(третья точка пересечения $L$ и $E$)}.
\]

\begin{center}\ChordPicture\end{center}

$(E(k),+)$ --- абелева группа. Абелевость очевидна.

Единица $O$ --- бесконечно удалённая точка:
\[
P+O=P\qquad\forall P\in E(k).
\]
Если $P=(x,y)$, то $-P=(x,-y)$.

\subsection*{Ассоциативность}
Достаточно доказать
\[
P_1(P_2+P_3)=(P_1+P_2)P_3.
\]
\begin{center}\AssociativityStart\end{center}

\newpage
% Рукопись, с. 4: основной рисунок.
\textbf{Ассоциативность.} Достаточно
\[
P_1(P_2+P_3)=(P_1+P_2)P_3.
\]

\begin{center}
\begin{tikzpicture}[x=1.48cm,y=.89cm,font=\small]
  \AssociativityConstruction{black}{P_1+P_3}
\end{tikzpicture}
\end{center}
\vspace{-2mm}
\[
C=\ell_1\ell_2\ell_3\quad\text{--- кубика},\qquad
\widetilde C=L_1L_2L_3.
\]
Они пересекаются в восьми точках с $E$. Следовательно, девятая точка
тоже лежит в пересечении, что и доказывает ассоциативность.

\subsection*{Вычислим координаты}
Пусть $P_1=(x_1,y_1)$, $P_2=(x_2,y_2)$. Найдём $P_1+P_2$.

\textbf{Случай 1: $P_1\ne P_2$.} Прямая через $P_1,P_2$:
\[
y=mx+b,\qquad m=\frac{y_2-y_1}{x_2-x_1}.
\]

\textbf{Случай 2: $P_1=P_2=P$.} Если $y=0$, то $2P=O$, то есть
$\ord(P)=2$. Если $y\ne0$, то
\[
y^2=x^3+Ax+B,\qquad 2yy'=3x^2+A,\qquad
y'=\frac{3x^2+A}{2y}.
\]

\textbf{Итог (подставить всё).}
\[
x_{P_1+P_2}=\left(\frac{y_2-y_1}{x_2-x_1}\right)^2-x_1-x_2,
\qquad P_1\ne P_2;
\]
\[
x_{2P}=\left(\frac{3x_1^2+A}{2y_1}\right)^2-2x_1.
\]

\newpage
\subsection*{Точки кручения}
\textbf{Замечание.} Пусть $P\in E(k)$. Тогда
\[
\ord(P)=2\Longleftrightarrow P=-P.
\]
То есть $(x,y)=(x,-y)$, откуда $y=0$, так как $\charop(k)\ne2$.

Пусть $k\subset\CC$ и
\[
y^2=x^3+Ax+B=(x-a_1)(x-a_2)(x-a_3).
\]
Тогда $\ord((a_i,0))=2$. Подгруппа точек порядка $2$:
\[
E[2]=
\begin{cases}
\Z/2\Z\times\Z/2\Z,&\text{если все }a_i\in k,\\
\Z/2\Z,&\text{если ровно один }a_i\in k,\\
0,&\text{иначе}.
\end{cases}
\]

% Рукопись, начало с. 5.
При $k\subseteq\CC$ записано
\[
E[N]=\Z/N\Z\times\Z/N\Z.
\]

$N$-кручение, точка порядка $N$.

\section{Кривые над конечным полем}
Рассмотрим случай кривых над $\F_p\simeq\Z/p\Z$.
Как найти $|E(\F_p)|$? В идеале --- найти $E(\F_p)$.

Для $E:\ y^2=x^3+Ax+B$:
\[
|E(\F_p)|=p+1+a_p,\qquad
a_p=\sum_{\alpha\in\F_p}\leg{\alpha^3+A\alpha+B}{p},
\]
где $\leg{\cdot}{p}$ --- символ Лежандра.

\begin{itemize}[leftmargin=1.6em]
\item Если $\leg{\alpha^3+A\alpha+B}{p}=1$, то $\alpha^3+A\alpha+B$
--- квадрат в $\F_p$; добавляются две точки в $E(\F_p)$.
\item Если $\leg{\alpha^3+A\alpha+B}{p}=-1$, то ничего не добавляется.
\item Если $\leg{\alpha^3+A\alpha+B}{p}=0$, то добавляется одна точка
$(\alpha,0)$.
\end{itemize}

То есть число точек равно
\begin{align*}
|E(\F_p)|
&=1+\sum_{\alpha\in\F_p}\left(\leg{\alpha^3+A\alpha+B}{p}+1\right)\\
&=(p+1)+\sum_{\alpha\in\F_p}\leg{\alpha^3+A\alpha+B}{p}.
\end{align*}

\textbf{Теорема (Хассе, 1933).}\quad $|a_p|\le2\sqrt p$.

\textbf{Замечание.} Следствие гипотез Вейля (Делинь, 1974).

\newpage
\subsection*{Формулы сложения}
Пусть $P_1=(x_1,y_1)$, $P_2=(x_2,y_2)$. Формулы для обеих координат:
\begin{align*}
x_{P_1+P_2}
&=\left(\frac{y_2-y_1}{x_2-x_1}\right)^2-x_1-x_2,\\
y_{P_1+P_2}
&=-\frac{y_2-y_1}{x_2-x_1}(x_{P_1+P_2}-x_1)-y_1;
\end{align*}
\begin{align*}
x_{2P}&=\left(\frac{3x_1^2+A}{2y_2}\right)^2-2x_1,\\
y_{2P}&=-\frac{3x_1^2+A}{2y_2}(x_{2P}-x_1)-y_1.
\end{align*}

\section{Пример над $\F_7$}
Рассмотрим
\[
E=\{y^2=x^3-2x-3\}.\qquad E(\F_7)\text{ --- ?}
\]
\[
\Delta=-4A^3-27B^2=4\cdot8-27\cdot9\ne0,
\]
поэтому это эллиптическая кривая.

% Рукопись, с. 6. Все промежуточные вычисления выполняются в F_7.
Возьмём $P=(3,2)$. Все дальнейшие вычисления --- в $\F_7$.

\textbf{$2P=(2,6)$.}
\[
x=\left(\frac{3\cdot9-2}{4}\right)^2-6
=\left(\frac{25}{4}\right)^2-6=2,
\]
\[
y=-(50)(2-3)-2=6.
\]

\textbf{$3P=(4,2)$.}
\[
(2,6)+(3,2)=(16-5=11=4,\;4\cdot2-6)=(4,2).
\]

\begin{align*}
4P&:\quad (4,2)+(3,2)=(0-7=0,\;5),\\
5P&:\quad (0,5)+(3,2)=(5,0),\\
6P&:\quad (5,0)+(3,2)=(0,2),\\
7P&:\quad (0,2)+(3,2)=(4,5).
\end{align*}

\textbf{$8P=2\cdot4P$.}
\[
\left(\frac{-2}{2\cdot5}\right)^2-0
=\left(\frac{-2}{3}\right)^2-0
\quad\Longrightarrow\quad 8P=(2,1).
\]
\[
9P:\quad (2,1)+(3,2)=(3,5),
\qquad 9P+P=O\quad\text{--- нуль группы}.
\]

\newpage
\textbf{Итог.}
\[
\ord(P)=10,\qquad \Z/10\Z\le E(\F_7).
\]
В рукописи записана оценка
\[
\bigl||E(\F_7)|-8\bigr|\le4.
\]

Если $\Z/10\Z\not\simeq\F_7$,
то
\[
\frac{|E(\F_7)|}{|\Z/10\Z|}\ge2,
\qquad\text{то есть }|E(\F_7)|\ge20,
\]
--- противоречие.

\section{Пример над $\F_{71}$}
\[
E=\{y^2=x^3-x\}.\qquad E(\F_{71})\text{ --- ?}
\]
Функция $f(x)=x^3-x$ нечётная:
\[
f(-\alpha)=-f(\alpha),\qquad f(x)=x(x-1)(x+1).
\]

$-1$ --- квадрат или нет в $\F_{71}$?
\[
(-1)^{(71-1)/2}=(-1)^{35}=-1,
\]
поэтому $-1$ --- квадратичный невычет. Тогда
\[
\leg{-\alpha}{71}
=\leg{-1}{71}\leg{\alpha}{71}
=-\leg{\alpha}{71}.
\]

При $p=71$:
\begin{align*}
\sum_{\alpha=1}^{35}\leg{\alpha^3-\alpha}{p}
+\sum_{\alpha=36}^{70}\leg{\alpha^3-\alpha}{p}
&=\sum_{\alpha=1}^{35}\leg{f(\alpha)}{p}
+\sum_{\alpha=1}^{35}\leg{f(-\alpha)}{p}\\
&=0.
\end{align*}
Следовательно,
\[
|E(\F_p)|=72.
\]

Пусть $P\in E(\F_{71})$. Тогда
\[
P=\underbrace{9P}_{\in E[8]}-\underbrace{8P}_{\in E[9]},
\qquad E(\F_{71})\simeq E[8]\times E[9].
\]

% Рукопись, с. 7.
\newpage
Каким может быть $E[8]$:
\[
\Z/8\Z,\qquad \Z/2\Z\times\Z/4\Z,\qquad
\Z/2\Z\times\Z/2\Z\times\Z/2\Z.
\]
Для $E[9]$:
\[
\Z/9\Z,\qquad \Z/3\Z\times\Z/3\Z.
\]

\subsection*{Подгруппа $E[8]$}
Пусть $P\in E[2]$, $P=(x_1,y_1)$. Тогда
\[
P=-P,\qquad (x_1,y_1)=(x_1,-y_1)\quad\Longrightarrow\quad y_1=0.
\]
\[
x^3-x=x(x-1)(x+1)\quad\text{--- три корня}.
\]
Элементы порядка $2$ в $\Z/2\Z\times\Z/4\Z$:
\[
(0,2),\quad(1,0),\quad(1,2)
\quad\Longrightarrow\quad E[8]\simeq\Z/2\Z\times\Z/4\Z.
\]

\subsection*{Подгруппа $E[9]$}
Пусть $P\in E[3]$. Тогда
\[
3P=O\quad\Longrightarrow\quad P=-2P
\quad\Longrightarrow\quad x_{2P}=x_P\pmod{71}.
\]
В рукописи получено
\[
3x^4-6x^2+1=0\pmod{71}.
\]

Корней не более четырёх; все они образуют пары $(r,-r)$.
То есть существует не более двух координат по $x$ для $3$-кручения,
а значит, не более четырёх точек в $E[3]$ (на аффинной части).
Следовательно,
\[
E[9]\simeq\Z/9\Z.
\]
То есть
\[
E(\F_{71})\simeq\Z/2\Z\times\Z/4\Z\times\Z/9\Z
\simeq\Z/2\Z\times\Z/36\Z.
\]
\end{document}
